\section*{第 \ref{ch:g7:geological-time} 章——地球的年龄}

\begin{solution}{exo:g7:geological-time:1}
(a) \omterm{def:g7:geological-time:unconformity}{不整合}；(b) \omterm{def:g7:geological-time:inclusion}{包裹体律}；(c) \omterm{def:g7:geological-time:superposition}{叠加原理}；(d) \omterm{def:g7:geological-time:fossil-succession}{化石
层序律}；(e) \omterm{def:g7:geological-time:cross-cutting}{切割律}。
\end{solution}

\begin{solution}{exo:g7:geological-time:2}
(a) 位于底部的红色砂岩最老，
位于顶部的黄色砂岩最年轻。(b) 所依据的是
\omterm{def:g7:geological-time:superposition}{叠加原理}。(c) 这些岩层自形成以来
未曾发生倒转。
\end{solution}

\begin{solution}{exo:g7:geological-time:3}
(a) 古生代、中生代、新生代。(b) 新生代。(c) 中生代。
(d) 侏罗纪紧接三叠纪之后；白垩纪
紧接古近纪之前。
\end{solution}

\begin{solution}{exo:g7:geological-time:4}
(a) 错误：\omterm{def:g7:geological-time:relative-dating}{相对定年}只能给\omterm{def:g3:stones-pebbles-sand:rock}{岩石}排序；以年为单位的
年龄是\omterm{def:g7:geological-time:relative-dating}{绝对年龄}。(b) 正确。(c) 错误：\omterm{def:g3:stones-pebbles-sand:pebble}{卵石}比包含它们的
砾岩更老（\omterm{def:g7:geological-time:inclusion}{包裹体律}）。(d) 正确
（\omterm{def:g7:geological-time:fossil-succession}{化石层序律}）。
\end{solution}

\begin{solution}{exo:g7:geological-time:5}
(a) 较年轻：它切穿了第 3 层，而切穿其他岩层的\omterm{def:g3:stones-pebbles-sand:rock}{岩石}较年轻
（\omterm{def:g7:geological-time:cross-cutting}{切割律}）。(b) 较老：第 4 层覆盖在它之上，并未被切穿。
(c) 第 1 层、第 2 层、第 3 层、花岗岩、第 4 层。(d) 花岗岩之上的岩层
被\omterm{def:g5:water-shapes-the-land:erosion}{侵蚀}掉，直至花岗岩出露地表；
随后第 4 层\omterm{def:g5:water-shapes-the-land:deposition}{沉积}在这一\omterm{def:g5:water-shapes-the-land:erosion}{侵蚀}面上，形成\omterm{def:g7:geological-time:unconformity}{不整合}。
\end{solution}

\begin{solution}{exo:g7:geological-time:6}
(a) $538.8 - 251.9 = 286.9$ 百万年。(b) $251.9 - 66 = 185.9$
百万年。(c) 6600 万年。(d) 古生代最长，
新生代最短（目前而言）。
\end{solution}

\begin{solution}{exo:g7:geological-time:7}
(a) $4540 - 538.8 = 4001.2$ 百万年。(b) $4001.2 \div 4540
\approx 0.881$：约 \textbf{88~\%}。(c) 古生代之前，有壳动物
尚不常见，无硬壳的生物几乎不会形成\omterm{def:g7:sediments-and-fossils:fossil}{化石}；
贝壳和骨骼易于石化，
只在古生代以后才常见。
\end{solution}

\begin{solution}{exo:g7:geological-time:8}
(a) $4540 \div 100 = 45.4$ 百万年每米。(b) $538.8 \div 45.4
\approx 11.9$：约 \qty{11.9}{m}，在终点线之前。(c) $66 \div
45.4 \approx 1.45$：约 \qty{1.45}{m}。(d) $0.3 \div 45.4 \approx
0.0066$~m，即约 \qty{0.66}{cm}，不到一厘米：
跑者几乎察觉不到。
\end{solution}

\begin{solution}{exo:g7:geological-time:9}
(a) 志留纪先形成。是的：较老的岩层位于
较年轻的岩层之下。(b) 灰色岩层水平\omterm{def:g5:water-shapes-the-land:deposition}{沉积}，随后被推至直立；
随后被\omterm{def:g5:water-shapes-the-land:erosion}{侵蚀}掉，红色砂层\omterm{def:g5:water-shapes-the-land:deposition}{沉积}在顶部。(c) 这是
两个年代相差悬殊的岩层之间的\omterm{def:g5:water-shapes-the-land:erosion}{侵蚀}面：漫长的时间
在此期间灰色岩层被抬升和\omterm{def:g5:water-shapes-the-land:erosion}{侵蚀}，
未留下任何岩层。
\end{solution}

\begin{solution}{exo:g7:geological-time:10}
(a) 两层含菊石的页岩年代相同，两层含海胆的石灰岩
年代也相同：\omterm{def:g7:sediments-and-fossils:fossil}{化石}相同即年代相同（依据
\omterm{def:g7:geological-time:fossil-succession}{化石层序律}）。(b) 最老的是悬崖 A 中含三叶虫的砂岩
（它位于含菊石的页岩之下，该页岩与悬崖 B 的
底层同龄）。最年轻的是悬崖
B 中含叶\omterm{def:g7:sediments-and-fossils:fossil}{化石}的黏土。(c) 自下而上：含三叶虫的砂岩；
含菊石的页岩；含海胆的石灰岩；含叶\omterm{def:g7:sediments-and-fossils:fossil}{化石}的黏土。
\end{solution}

\begin{solution}{exo:g7:geological-time:11}
(a) 片岩。它位于砾岩之下（\omterm{def:g7:geological-time:superposition}{叠加原理}），砾岩中的
\omterm{def:g3:stones-pebbles-sand:pebble}{卵石}是片岩碎屑（\omterm{def:g7:geological-time:inclusion}{包裹体律}）；
\omterm{def:g5:water-shapes-the-land:erosion}{侵蚀}面也切穿了片岩。(b) 较晚发生：砂岩先水平\omterm{def:g5:water-shapes-the-land:deposition}{沉积}，
之后才被倾斜。(c) 较早发生：黏土覆盖在\omterm{def:g7:earthquakes:fault}{断层}上，
未被错断。(d) 片岩形成，随后被\omterm{def:g5:water-shapes-the-land:erosion}{侵蚀}成起伏面；
随后依次\omterm{def:g5:water-shapes-the-land:deposition}{沉积}砾岩、砂岩和石灰岩；
岩层被倾斜，\omterm{def:g7:earthquakes:fault}{断层}将其错断——
仅凭该剖面无法判断这两者的先后；
倾斜岩层和\omterm{def:g7:earthquakes:fault}{断层}被削平，黏土\omterm{def:g5:water-shapes-the-land:deposition}{沉积}在顶部。
(e) 两处：片岩与砾岩之间的起伏面，
以及倾斜岩层与平坦黏土层之间的\omterm{def:g5:water-shapes-the-land:erosion}{侵蚀}面。
\end{solution}

\begin{solution}{exo:g7:geological-time:12}
(a) 365 天的一半是 182.5 天。1 月至 6 月共
$31+28+31+30+31+30 = 181$ 天，因此中点落在 7 月初之后 1.5 天，
即 \textbf{7 月 2 日正午}。(b) $4540 \div 2 = 2270$ 百万
年。(c) $80 \div \num{4540000000} \times \num{31536000} \approx 0.56$：
约 \textbf{0.6~秒}。
\end{solution}

\begin{solution}{pb:g7:geological-time:1}
\textbf{1.} 45.4 亿年。

\textbf{2.} $4540 \div 365 \approx 12.4$ 百万年。

\textbf{3.} $\num{4540000000} \div (365 \times 24) = \num{4540000000}
\div 8760 \approx \num{518000}$ 年（用取整后的 \num{12.4}
百万，$\num{12400000} \div 24 \approx \num{517000}$）。

\textbf{4.} $\num{4540000000} \div \num{525600} \approx 8638$：约
\num{8600} 年。

\textbf{5.} $4540 - 4030 = 510$ 百万年（地球形成之后）；
$510 \div 12.4 \approx 41$ 天。1 月有 31 天，因此 41 天
至 \textbf{2 月 11 日}初已过去。

\textbf{6.} $4540 - 3500 = 1040$ 百万年；$1040 \div 12.4 \approx
83.9$ 天（用未取整的 $12.44$ 时为 83.6 天）。1 月和 2 月共
59 天，进入 3 月后 24.6 至 24.9 天，即 \textbf{3 月 25 日}。

\textbf{7.} $66 \div 12.4 \approx 5.3$ 天，相当于年末前
即 1 月 1 日后 $365 - 5.3 = 359.7$ 天。前十一个月
共 334 天，进入 12 月后 25.7 天，即 \textbf{12 月 26 日}。

\textbf{8.} $538.8 \div 4540 \approx 0.119$：\textbf{11.9~\%} 的年份
尚未过去，即 $0.119 \times 365 \approx 43$ 天。$365 - 43 =
322$ 天，即 1 月 1 日后 322 天，开始于 \textbf{11 月}（
11 月 18 日）。

\textbf{9.} $\num{10000} \div 8600 \approx 1.2$ 分钟。

\textbf{10.} $5000 \div 8600 \approx 0.58$ 分钟，而 $0.58 \times 60
\approx 35$：约在午夜前 \textbf{35 秒}。

\textbf{11.} $\num{300000} \div 8600 \approx 34.9$：约在午夜前 \textbf{35
分钟}。

\textbf{12.} 12 月 31 日午夜前 35 分钟：人类
约出现在 \textbf{23:25}，处于全年最后一小时内。
\end{solution}
